\section{Descent and minimality}\label{dp:sec:descent}

We prove that divided-power analyticity descends along isometries of the target and along surjections of the source. Then we write each divided-power germ as an analytic germ after one enlargement of the target and one cover of the source. It follows that the divided-power germs are the smallest class containing the analytic germs with these two descents (Theorem~\ref{dp:thm:minimal}). Between ultrametric spaces source descent alone suffices (Theorem~\ref{dp:thm:ultrametric-minimal}). The section ends with what this minimality depends on (Example~\ref{dp:ex:l1}).

The analytic and the AA germs are as in \cite[Definition~14.7]{I}, with $M$ for the bound. Recall $C_1=8/\abs\varpi^2$ and $C_0=\mathrm eC_1$.

\subsection{The two descents}\label{dp:ssec:two-descents}

\begin{theorem}[target descent]\label{dp:thm:target-descent}
Let $j\colon W\to Z$ be a linear isometry of normed spaces, $U\subseteq P$ open, $f\colon U\to W$ and $x_0\in U$. Suppose that $j\circ f$ is DP at $x_0$ with radius $R$, bound $M$ and coefficient laws $a'_n\colon P\to Z$.
\begin{enumerate}
\item If every $a'_n$ takes values in $j(W)$, then $f$ is DP at $x_0$ with radius $R$ and bound $M$, with coefficient laws $a_n=j^{-1}\circ a'_n$ and $\norm{a_n}=\norm{a'_n}$.
\item If $j(W)$ is closed in $Z$, for instance if $W$ is complete, the hypothesis of~(1) holds. So divided-power analyticity, at a point or on an open set, descends along linear isometries with closed range, with the same radius and bound.
\end{enumerate}
\end{theorem}

\begin{proof}
\emph{Step 1 (item (1)).} Let $h_1,\dots,h_r\in P$ and $\abs\alpha=n$. The vector $(a'_n)^{[\alpha]}(h_1,\dots,h_r)$ is a coefficient of the polynomial $\lambda\mapsto a'_n(\sum_i\lambda_ih_i)$ on $K^r$, whose values lie in $j(W)$. As $K$ is infinite, interpolation shows that this coefficient lies in $j(W)$. So $a_n(\sum_i\lambda_ih_i)=\sum_\alpha\lambda^\alpha j^{-1}(a'_n)^{[\alpha]}(h_1,\dots,h_r)$ is a homogeneous polynomial with coefficients in $W$, and $\norm{a_n}=\norm{a'_n}\le MR^{-n}$ because $j$ is isometric. For $\norm h<R$ the partial sums $\sum_{n\le N}a_n(h)\in W$ have images converging to $jf(x_0+h)$, so they converge to $f(x_0+h)$.

\emph{Step 2 (item (2)).} For $h=0$ the values $a'_n(0)$ are $jf(x_0)$ and $0$. Fix $h\ne0$ and put $\rho=R/\norm h$. For $\abs s<\rho$ we have $\psi(s)=jf(x_0+sh)=\sum_ns^na'_n(h)$, with $\norm{a'_n(h)}\rho^n\le M$. We show $a'_n(h)\in j(W)$ by induction on $n$. First, $a'_0(h)=jf(x_0)$. If $a'_m(h)\in j(W)$ for $m<n$, then for $0<\abs s<\rho$,
\begin{gather*}
 \delta_n(s)=s^{-n}\Bigl(\psi(s)-\sum_{m<n}s^ma'_m(h)\Bigr)\in j(W),\\
 \norm{\delta_n(s)-a'_n(h)}\le\sum_{m>n}M\rho^{-m}\abs s^{m-n}=M\rho^{-n}\frac{\abs s/\rho}{1-\abs s/\rho}.
\end{gather*}
As $K$ is nontrivially valued, $s$ can tend to $0$, so $a'_n(h)$ lies in the closure of $j(W)$, which is $j(W)$. If $W$ is complete, then $j(W)$ is complete, hence closed.
\end{proof}

Without closed range, target descent can fail at every point.

\begin{example}[a dense isometry]\label{dp:ex:dense-isometry}
Let $K$ be complete. Let $Z=c_0(\N,K)$ and $\Delta=\{y\in K:\abs y<1\}$. Put $\mathsf F(y)=(y^i)_{i\ge0}\in Z$, with $0^0=1$. Let $W$ be the linear span of $\mathsf F(\Delta)$, with the norm of $Z$, and let $j\colon W\to Z$ be the inclusion. Then:
\begin{enumerate}
\item $j\circ\mathsf F$ is analytic at every point of $\Delta$, with radius $1$ and bound $1$. So is $\mathsf F$ as a map into the closure $\overline W$ of $W$ in $Z$, which is a completion of $W$.
\item $\mathsf F\colon\Delta\to W$ is DP at no point of $\Delta$, and it is differentiable at no point as a $W$-valued map.
\end{enumerate}
So in \cite[Definition~14.7]{I} the condition on diagonals cannot be dropped: the maps $\beta_n$ of the proof expand $j\circ\mathsf F$ with no condition on their diagonals, while $\mathsf F$ is AA at no point, as AA germs are DP \cite[Proposition~14.9(1)]{I}.
\end{example}

\begin{proof}
\emph{Step 1 (item (1)).} Let $y_0\in\Delta$. For $\abs s<1$ we have $\abs{y_0+s}<1$ and, coordinatewise, $(y_0+s)^i=\sum_{n\le i}\binom iny_0^{i-n}s^n$. Put $z_n=(\binom iny_0^{i-n})_{i\ge0}$, with entry $0$ for $i<n$. Then $z_n\in Z$ and $\norm{z_n}\le1$, as $\abs{\binom in}\le1$ and $\abs{y_0}^{i-n}\to0$ as $i\to\infty$. Since $Z$ is complete, $\sum_ns^nz_n$ converges in $Z$, and coordinatewise it converges to $\mathsf F(y_0+s)$. So $j\mathsf F(y_0+s)=\sum_n\beta_n(s,\dots,s)$ with the $n$-linear maps $\beta_n(s_1,\dots,s_n)=s_1\cdots s_nz_n$, of norm at most $1$. Now apply Theorem~\ref{dp:thm:target-descent}(2) to the inclusion of the closed subspace $\overline W$ in $Z$, with Lemma~\ref{dp:lem:coefficients-unique}(2). This gives $z_n=\beta_n(1,\dots,1)\in\overline W$, so the $\beta_n$ expand $\mathsf F$ as a map into $\overline W$.

\emph{Step 2 (the vector that leaves $W$).} Suppose $\mathsf F$ is DP into $W$ at $y_0$, with laws $a_n$. By Lemma~\ref{dp:lem:coefficients-unique}(2), $ja_n(s)=\beta_n(s,\dots,s)$. So $z_1=(\mathsf d_i)_{i\ge0}$, with $\mathsf d_i=iy_0^{i-1}$ and $\mathsf d_0=0$, lies in $W$. The vector $z_1$ is also the limit in $Z$ of the difference quotients $s^{-1}(\mathsf F(y_0+s)-\mathsf F(y_0))$, so a derivative of $\mathsf F$ in $W$ at $y_0$ would also force $z_1\in W$.

\emph{Step 3 ($z_1\notin W$).} Let $\sh$ be the shift, $(\sh\mathsf t)_i=\mathsf t_{i+1}$, on sequences $\mathsf t$. The identity $(i+2)y^{i+1}-2y(i+1)y^i+y^2iy^{i-1}=0$ gives $(\sh-y_0)^2z_1=0$, while $((\sh-y_0)z_1)_0=\mathsf d_1-y_0\mathsf d_0=1$. Suppose $z_1=\sum_{l\le r}\lambda_l\mathsf F(y_l)$ with distinct $y_l\in\Delta$. Since $(\sh-y_l)\mathsf F(y_l)=0$, also for $y_l=0$, we get $\prod_l(\sh-y_l)z_1=0$. The polynomials $\mathsf g\in K[X]$ with $\mathsf g(\sh)z_1=0$ form an ideal, and its monic generator divides both $(X-y_0)^2$ and the squarefree $\prod_l(X-y_l)$. So it divides $X-y_0$, and $(\sh-y_0)z_1=0$. This is a contradiction.
\end{proof}

\begin{corollary}[values in the completion]\label{dp:cor:completion}
Let $U\subseteq P$ be open, $f\colon U\to W$ and $x_0\in U$. Then $f$ is DP at $x_0$ as a $W$-valued map if and only if it is DP at $x_0$ as a $\widehat W$-valued map and, for every $h\in P$ and $n\ge0$, the coefficient of $s^n$ in $s\mapsto f(x_0+sh)$, computed in $\widehat W$, lies in $W$. Radii and bounds agree. The second condition can fail at every point of an open set on which $f$ is analytic into $\widehat W$ (Example~\ref{dp:ex:dense-isometry}).
\end{corollary}

\begin{proof}
Apply Theorem~\ref{dp:thm:target-descent}(1) with $j\colon W\to\widehat W$, and Lemma~\ref{dp:lem:coefficients-unique}(1). The converse is immediate.
\end{proof}

Recall from Lemma~\ref{dp:lem:strong-formal}(2) that a bounded linear surjection $Q\colon P'\to P$ has \emph{preimage constant} $\mathsf c>0$ if every $h\in P$ has a preimage $\tilde h$ with $\norm{\tilde h}\le\mathsf c\norm h$.

\begin{proposition}[source descent]\label{dp:prop:source-descent}
Let $Q\colon P'\to P$ be a bounded linear surjection with preimage constant $\mathsf c$. Let $U\subseteq P$ be open, $f\colon U\to W$, $\tilde x_0\in Q^{-1}(U)$ and $x_0=Q(\tilde x_0)$.
\begin{enumerate}
\item If $f\circ Q$ is DP at $\tilde x_0$ with radius $R'$, bound $M$ and coefficient laws $b_n$, then $b_n=a_n\circ Q$ for bounded laws $a_n\colon P\to W$ with $\norm{a_n}\le\mathsf c^n\norm{b_n}$, and $f$ is DP at $x_0$ with radius $R'/\mathsf c$, bound $M$ and coefficient laws $a_n$.
\item If $f$ is DP at $x_0$ with radius $R$, bound $M$ and coefficient laws $a_n$, then $f\circ Q$ is DP at $\tilde x_0$ with radius $R/\norm Q$ (any radius if $Q=0$), bound $M$ and coefficient laws $a_n\circ Q$.
\end{enumerate}
\end{proposition}

\begin{proof}
\emph{Step 1 ($b_n$ is constant on fibers).} Let $\tilde h\in P'$ and $\mathsf k\in\ker Q$. For small $s$ we have $f\circ Q(\tilde x_0+s(\tilde h+\mathsf k))=f\circ Q(\tilde x_0+s\tilde h)$. Both sides are power series in $s$ with geometric bounds, with coefficients $b_n(\tilde h+\mathsf k)$ and $b_n(\tilde h)$. By Lemma~\ref{dp:lem:coefficients-unique}(1), $b_n(\tilde h+\mathsf k)=b_n(\tilde h)$. So $a_n(h)=b_n(\tilde h)$, for any preimage $\tilde h$ of $h$, is well defined, and $b_n=a_n\circ Q$.

\emph{Step 2 ($a_n$ is a bounded law).} For $h_1,\dots,h_r\in P$ with preimages $\tilde h_i$, the vector $\sum_i\lambda_i\tilde h_i$ is a preimage of $\sum_i\lambda_ih_i$. So $a_n(\sum_i\lambda_ih_i)=b_n(\sum_i\lambda_i\tilde h_i)$ is a homogeneous polynomial of degree $n$ with coefficients in $W$. Taking $\norm{\tilde h}\le\mathsf c\norm h$ gives $\norm{a_n(h)}\le\norm{b_n}\mathsf c^n\norm h^n$.

\emph{Step 3 (expansion).} For $\norm h<R'/\mathsf c$ take $\tilde h$ with $\norm{\tilde h}\le\mathsf c\norm h<R'$. Then $x_0+h\in U$ and $f(x_0+h)=f\circ Q(\tilde x_0+\tilde h)=\sum_nb_n(\tilde h)=\sum_na_n(h)$. Also $\norm{a_n}\le\mathsf c^nMR'^{-n}=M(R'/\mathsf c)^{-n}$.

\emph{Step 4 (item (2)).} The map $a_n\circ Q$ is a bounded law with $\norm{a_n\circ Q}\le\norm Q^n\norm{a_n}$, and $f\circ Q(\tilde x_0+\tilde h)=\sum_na_n(Q\tilde h)$ for $\norm Q\norm{\tilde h}<R$.
\end{proof}

\subsection{The canonical cover}\label{dp:ssec:cover}

\begin{definition}[canonical cover]\label{dp:def:cover}
Recall $\ell^1(J)$ and $\ell^1_{00}(J)$ from Section~\ref{dp:sec:conv}, and write $e_\theta$ for the unit vectors. For a normed space $P$ let $I_P=\{h\in P:\abs\varpi<\norm h\le1\}$, and let $q_P\colon\ell^1_{00}(I_P)\to P$ be the linear map with $q_P(e_h)=h$. If $K$ and $P$ are complete, let $\hat q_P\colon\ell^1(I_P)\to P$, $\mathsf v\mapsto\sum_{h\in I_P}\mathsf v_hh$.
\end{definition}

The maps $q_P$, $\hat q_P$ are contractive, as $\norm{\sum_h\mathsf v_hh}\le\sum_h\abs{\mathsf v_h}\norm h\le\norm{\mathsf v}$. Each has preimage constant $\abs\varpi^{-1}$. Indeed, for $h\ne0$ choose $\lambda\in K$ with $\abs\varpi<\abs\lambda\norm h\le1$. Then $\lambda h\in I_P$, $h=q_P(\lambda^{-1}e_{\lambda h})$ and $\abs\lambda^{-1}<\abs\varpi^{-1}\norm h$.

For a set $J$ with a fixed total order, and a law $b\colon\ell^1_{00}(J)\to W$ of degree $n$, we use the following construction. For a nondecreasing word $\mathsf w=(\mathsf w_1\le\dots\le\mathsf w_n)$ in $J$, with distinct letters $\theta_1<\dots<\theta_r$ occurring $\alpha_1,\dots,\alpha_r$ times, put $v_{\mathsf w}=b^{[\alpha]}(e_{\theta_1},\dots,e_{\theta_r})$. For finitely supported $h^1,\dots,h^n$ put
\begin{equation}\label{dp:eq:word-lift}
 \beta_b(h^1,\dots,h^n)=\sum_{\mathsf w}h^1_{\mathsf w_1}\cdots h^n_{\mathsf w_n}\,v_{\mathsf w},
\end{equation}
a finite sum over nondecreasing words. This map is $n$-linear. Let $h$ be supported in a finite set $J_0\subseteq J$. By Lemma~\ref{dp:lem:elementary}(3), $b(h)$ is the sum over $\alpha\in\N^{J_0}$ with $\abs\alpha=n$ of $h^\alpha b^{[\alpha]}((e_\theta)_{\theta\in J_0})$, and arguments with $\alpha_\theta=0$ can be deleted. Each monomial $h^\alpha$ comes from exactly one nondecreasing word. So $\beta_b(h,\dots,h)=b(h)$. For $n=0$, $\beta_b$ is the constant $b$.

\begin{lemma}[laws on $\ell^1_{00}$ are diagonals]\label{dp:lem:laws-cover}
Let $W$ be normed, and fix a total order on $J$. Every bounded law $b\colon\ell^1_{00}(J)\to W$ of degree $n$ is the diagonal of the $n$-linear map $\beta_b\colon\ell^1_{00}(J)^n\to W$, with $\norm{\beta_b}\le C_0^n\,n!\,\norm b$. So $b$ is analytic everywhere.
\end{lemma}

\begin{proof}
By Theorem~\ref{dp:thm:mixed}(3), $\norm{v_{\mathsf w}}\le C_0^n\norm b\,n!/\prod_c\alpha_c!\le C_0^nn!\norm b$. Also $\sum_{\mathsf w}\prod_c\abs{h^c_{\mathsf w_c}}\le\prod_c\norm{h^c}$. The last assertion is \cite[Proposition~2.5]{I}, with $j=\id_W$.
\end{proof}

The factor $n!$ is not geometric. Removing it is what the enlargement in Theorem~\ref{dp:thm:normal-form} is for. Recall $\operatorname{lift}(a)$ from Definition~\ref{dp:def:strong-norm}(1).

\begin{lemma}[lifts after the cover]\label{dp:lem:lift-cover}
Let $W$ be normed and let $a\colon P\to W$ be a law of degree $n$. Fix a total order on $I_P$. The map $\beta_{a\circ q_P}$ of~\eqref{dp:eq:word-lift}, for $J=I_P$, has norm at most $\lvert\!\lvert\!\lvert a\rvert\!\rvert\!\rvert$. So $\operatorname{lift}(a\circ q_P)\le\lvert\!\lvert\!\lvert a\rvert\!\rvert\!\rvert$.
\end{lemma}

\begin{proof}
Let $b=a\circ q_P$, and take $\beta_b$ of~\eqref{dp:eq:word-lift} for $J=I_P$. For a word with letters $\xi_1<\dots<\xi_r$ in $I_P$ and multiplicities $\alpha$ we have $v_{\mathsf w}=a^{[\alpha]}(\xi_1,\dots,\xi_r)$, as $q_P$ is linear. As $\norm{\xi_c}\le1$, Definition~\ref{dp:def:strong-norm} gives $\norm{v_{\mathsf w}}\le\lvert\!\lvert\!\lvert a\rvert\!\rvert\!\rvert$. Conclude as in Lemma~\ref{dp:lem:laws-cover}.
\end{proof}

\begin{remark}[failures of the two descents for analyticity]\label{dp:rem:an-no-descent}
\emph{Target.} Let $\operatorname{char}K=p$, $k\ge p$, and let $E,F$ be as in \cite[Theorem~7.2(1)]{I}. Then $(j_{E,F}\circ{-})\circ A^k_{E,E;F}$ is analytic everywhere, $j_{E,F}\circ{-}$ has closed range, and $A^k_{E,E;F}$ is analytic at no point \cite[Remark~14.13]{I}. In every characteristic, the germ $\mathsf f$ of Example~\ref{dp:ex:l1} shows the same.

\emph{Source.} Let $\operatorname{char}K=p$ and let $a$ be the law of Proposition~\ref{dp:prop:degree-p}(2) with $n=p$, into the completion of $\Gamma^p(\mathcal E)$. By Lemma~\ref{dp:lem:laws-cover}, $a\circ q_{K\oplus\mathcal E}$ is analytic everywhere. By \cite[Proposition~14.9(2)]{I} with $j=\id$, $a$ is analytic at no point, since it is not ambient.
\end{remark}

\begin{theorem}[normal form]\label{dp:thm:normal-form}
Let $W$ be a Banach space, $U\subseteq P$ open, $f\colon U\to W$, $x_0\in U$, and $\tilde x_0\in\ell^1_{00}(I_P)$ with $q_P(\tilde x_0)=x_0$. Let $\mathfrak U$, $u_0$ be as in \cite[(14.8)]{I}, and let $j_W\colon W\to W\hot_\pi\mathfrak U$, $w\mapsto w\otimes u_0$, which is a linear isometry into a Banach space \cite[Remark~14.41]{I}. The following are equivalent:
\begin{enumerate}
\item $f$ is DP at $x_0$;
\item $f\circ q_P$ is AA at $\tilde x_0$;
\item there are a Banach space $Z$ and a linear isometry $j\colon W\to Z$ such that $j\circ f\circ q_P$ is analytic at $\tilde x_0$;
\item $j_W\circ f\circ q_P$ is analytic at $\tilde x_0$.
\end{enumerate}
Radius $R$ and bound $M$ in~(1) give radius $R/C_0$ and bound $M$ in~(2), (3), (4). Radius $R'$ and bound $M$ in~(3) or~(4) give radius $\abs\varpi R'$ and bound $M$ in~(1).
\end{theorem}

\begin{proof}
\emph{Step 1 ((1) implies (2)).} By Proposition~\ref{dp:prop:source-descent}(2) and $\norm{q_P}\le1$, $f\circ q_P$ is DP at $\tilde x_0$ with radius $R$, bound $M$ and laws $b_n=a_n\circ q_P$. By \cite[Theorem~14.39(2)]{I} and $\norm{L_{b_n}(t)}\le\norm{b_n}\norm t_\Gamma$ (Theorem~\ref{dp:thm:normed-roby}(1)), we get
\begin{equation}\label{dp:eq:amb-cover}
 \operatorname{amb}(b_n)\le C_0^nMR^{-n}=M(R/C_0)^{-n},
\end{equation}
and for $n=0$ the law is a constant of norm at most $M$. As $C_0>1$, \cite[Proposition~14.9(1)]{I} gives that $f\circ q_P$ is AA with radius $R/C_0$ and bound $M$.

\emph{Step 2 ((2) implies (3)).} This is \cite[Remark~14.13]{I}, with the same radius and bound.

\emph{Step 3 ((3) implies (1)).} The map $j\circ f\circ q_P$ is DP at $\tilde x_0$ with radius $R'$ and bound $M$ \cite[Proposition~14.9(1)]{I}. As $W$ is complete, Theorem~\ref{dp:thm:target-descent}(2) gives the same for $f\circ q_P$. Proposition~\ref{dp:prop:source-descent}(1) with $\mathsf c=\abs\varpi^{-1}$ gives that $f$ is DP at $x_0$ with radius $\abs\varpi R'$ and bound $M$.

\emph{Step 4 ((1) implies (4)).} This is the construction of \cite[Remark~14.41]{I}, stated there on $\ell^1(J)$. On $\ell^1_{00}(J)$, with $J=I_P$, every sum below is finite. Recall that $u_0=\sum_{i\le N}\mathsf u_{N,i}$ with $\norm{\mathsf u_{N,i}}\le1/N$. Fix $n\ge1$. For finitely supported $\alpha\in\N^J$ with $\abs\alpha=n$ let $O_\alpha(e)$ be the orbit sum of the unit vectors in the support of $\alpha$, and $w_\alpha=L_{b_n}(O_\alpha(e))$. There are $N_\alpha=n!/\prod_\theta\alpha_\theta!$ words $\mathsf w\in J^n$ of content $\alpha$, that is, in which each $\theta$ occurs $\alpha_\theta$ times. So $\pi(O_\alpha(e))\le N_\alpha$ and $\norm{w_\alpha}\le\operatorname{amb}(b_n)N_\alpha$. Number these words, and put $\mathsf z_{\mathsf w}=w_\alpha\otimes\mathsf u_{N_\alpha,i}$ if $\mathsf w$ is the $i$-th. Then $\norm{\mathsf z_{\mathsf w}}\le\operatorname{amb}(b_n)$, and the $\mathsf z_{\mathsf w}$ of content $\alpha$ sum to $w_\alpha\otimes u_0$. Put
\[
 \beta_n(h^1,\dots,h^n)=\sum_{\mathsf w\in J^n}h^1_{\mathsf w_1}\cdots h^n_{\mathsf w_n}\,\mathsf z_{\mathsf w}.
\]
Then $\norm{\beta_n}\le\operatorname{amb}(b_n)\le M(R/C_0)^{-n}$ by~\eqref{dp:eq:amb-cover}, and by~\eqref{dp:eq:orbit},
\[
 \beta_n(h,\dots,h)=\sum_\alpha h^\alpha w_\alpha\otimes u_0=L_{b_n}(h^{\otimes n})\otimes u_0=j_W(b_n(h)).
\]
With $\beta_0=j_W(f(x_0))$ this expands $j_W\circ f\circ q_P$ at $\tilde x_0$ with radius $R/C_0$ and bound $M$.

\emph{Step 5 ((4) implies (3)).} Take $j=j_W$.
\end{proof}

\subsection{Minimality}\label{dp:ssec:minimal}

Here a \emph{germ} is a quadruple $(P,U,f,x_0)$: $P$ a normed space, $U\subseteq P$ open, $x_0\in U$ and $f\colon U\to W$ with $W$ a Banach space. Consider three properties of a class $\cC$ of germs.
\begin{itemize}
\item[(C1)] Every germ $(P,U,f,x_0)$ with $f$ analytic at $x_0$ is in $\cC$.
\item[(C2)] \emph{(Target descent.)} If $j\colon W\to Z$ is a linear isometry of Banach spaces and $(P,U,j\circ f,x_0)\in\cC$, then $(P,U,f,x_0)\in\cC$.
\item[(C3)] \emph{(Source descent.)} If $Q\colon P'\to P$ is a bounded linear surjection with a preimage constant, $\tilde x_0\in Q^{-1}(U)$ and $(P',Q^{-1}(U),f\circ Q,\tilde x_0)\in\cC$, then $(P,U,f,Q(\tilde x_0))\in\cC$.
\end{itemize}

\begin{theorem}[the smallest class with these descents]\label{dp:thm:minimal}
\leavevmode
\begin{enumerate}
\item The DP germs satisfy (C1), (C2), (C3).
\item Every class of germs with (C1), (C2), (C3) contains every DP germ. It suffices to require (C1) for the sources $\ell^1_{00}(J)$, (C2) for the isometries $j_W$ of Theorem~\ref{dp:thm:normal-form} and (C3) for the canonical covers $q_P$.
\item So the DP germs are the smallest class of germs with (C1), (C2), (C3).
\item Let $K$ be complete. The DP germs with Banach sources are the smallest class of germs with Banach sources that has (C1) and (C2) and has (C3) for surjections between Banach spaces.
\end{enumerate}
\end{theorem}

\begin{proof}
\emph{Step 1 (item (1)).} (C1) is Theorem~\ref{dp:thm:calculus}(3). For (C2), an isometry from a Banach space has closed range, and Theorem~\ref{dp:thm:target-descent}(2) applies. (C3) is Proposition~\ref{dp:prop:source-descent}(1).

\emph{Step 2 (item (2)).} Let $(P,U,f,x_0)$ be a DP germ, and choose $\tilde x_0$ with $q_P(\tilde x_0)=x_0$. By Theorem~\ref{dp:thm:normal-form}(4), $(\ell^1_{00}(I_P),q_P^{-1}(U),j_W\circ f\circ q_P,\tilde x_0)\in\cC$ by (C1). So $(\ell^1_{00}(I_P),q_P^{-1}(U),f\circ q_P,\tilde x_0)\in\cC$ by (C2), and $(P,U,f,x_0)\in\cC$ by (C3). Item~(3) follows from items~(1) and~(2).

\emph{Step 3 (item (4)).} Let $P$ be a Banach space and choose $\tilde x_0\in\ell^1(I_P)$ with $\hat q_P(\tilde x_0)=x_0$. As $K$ is complete, $\ell^1(I_P)$ is a Banach space, and $\hat q_P$ is contractive with preimage constant $\abs\varpi^{-1}$. By \cite[Theorem~14.39(2)]{I} for $\ell^1(I_P)$, Steps~1 and~2 of the proof of Theorem~\ref{dp:thm:normal-form} apply with $\hat q_P$ in place of $q_P$. They give a Banach space $Z$ and a linear isometry $j\colon W\to Z$ with $j\circ f\circ\hat q_P$ analytic at $\tilde x_0$. Then (C1), (C2), (C3) apply as in Step~2, and all sources involved are Banach spaces.
\end{proof}

\begin{proposition}[the two other classes]\label{dp:prop:other-classes}
\leavevmode
\begin{enumerate}
\item The analytic germs satisfy (C1). Over every $K$ they fail (C2). If $\operatorname{char}K=p$ they fail (C3). If $\operatorname{char}K=0$, then (C2) and (C3) hold for germs whose source $P$ is ultrametric, the source $P'$ of $Q$ in (C3) being arbitrary.
\item The AA germs are the smallest class of germs with (C1) and (C2). If $\operatorname{char}K=p$ they fail (C3).
\end{enumerate}
\end{proposition}

\begin{proof}
\emph{Step 1 (item (1)).} The failures are Example~\ref{dp:ex:l1} and Remark~\ref{dp:rem:an-no-descent}. Let $\operatorname{char}K=0$ and $P$ be ultrametric. If $j\circ f$ is analytic, it is DP, so $f$ is DP (Theorem~\ref{dp:thm:target-descent}(2)), so $f$ is analytic \cite[Proposition~14.40(2)]{I}. If $f\circ Q$ is analytic, it is DP, so $f$ is DP (Proposition~\ref{dp:prop:source-descent}(1)), and again $f$ is analytic.

\emph{Step 2 (item (2): minimality).} Analytic germs are AA, with $j=\id$. If $\cC$ has (C1) and (C2), and $f$ is AA at $x_0$, then \cite[Remark~14.13]{I} gives a Banach space $Z$ and an isometry $j$ with $j\circ f$ analytic. So $f\in\cC$.

\emph{Step 3 (item (2): the AA germs have (C2)).} Let $j\colon W\to Z$ be an isometry of Banach spaces with $j\circ f$ AA. By \cite[Remark~14.13]{I} there is an isometry $j'$ of $Z$ into a Banach space with $j'\circ j\circ f$ analytic. So $f$ is AA, through $j'\circ j$, by the same remark.

\emph{Step 4 (item (2): failure of (C3)).} Let $a$ be the law of Remark~\ref{dp:rem:an-no-descent}. It has Banach values and is not ambient. If it were AA at a point, \cite[Remark~14.13]{I} and \cite[Proposition~14.9(2)]{I} would make it ambient. And $a\circ q_{K\oplus\mathcal E}$ is analytic, hence AA, everywhere.
\end{proof}

In characteristic $0$, on non-ultrametric sources, the analytic germs with ultrametric Banach targets have (C3) (Corollary~\ref{dp:cor:two-gaps}). For non-ultrametric sources and targets in characteristic $0$, whether they have (C3) is open.

The descents (C2), (C3) are properties of $C^\infty$ maps. For (C2) this is \cite[Lemma~2.4]{I}. For (C3):

\begin{proposition}[$C^m$ descends along source quotients]\label{dp:prop:cn-descends}
Let $Q\colon P'\to P$ be a bounded linear surjection with preimage constant $\mathsf c$, $U\subseteq P$ open, $W$ normed, $m\in\N\cup\{\infty\}$ and $f\colon U\to W$. If $f\circ Q$ is $C^m$ on $Q^{-1}(U)$, then $f$ is $C^m$ on $U$.
\end{proposition}

\begin{proof}
\emph{Step 1 ($m=0$).} Let $x_0=Q(\tilde x_0)\in U$. For $x$ near $x_0$ choose $\tilde x$ with $Q\tilde x=x$ and $\norm{\tilde x-\tilde x_0}\le\mathsf c\norm{x-x_0}$. Then $f(x)-f(x_0)=f\circ Q(\tilde x)-f\circ Q(\tilde x_0)\to0$.

\emph{Step 2 ($m=1$).} Let $x=Q\tilde x$ and $T=D(f\circ Q)(\tilde x)$. For $\mathsf k\in\ker Q$ and small $s$ we have $f\circ Q(\tilde x+s\mathsf k)=f(x)$, so $T\mathsf k=0$. Hence $T=T_0\circ Q$ for a unique linear $T_0\colon P\to W$, and $\norm{T_0h}=\norm{T\tilde h}\le\mathsf c\norm T\norm h$. Also $\norm{f(x+h)-f(x)-T_0h}=\norm{f\circ Q(\tilde x+\tilde h)-f\circ Q(\tilde x)-T\tilde h}=o(\norm{\tilde h})=o(\norm h)$, so $Df(x)=T_0$. For $x'$ near $x$ choose $\tilde x'$ with $Q\tilde x'=x'$ and $\norm{\tilde x'-\tilde x}\le\mathsf c\norm{x'-x}$. Then $\norm{Df(x')-Df(x)}\le\mathsf c\norm{D(f\circ Q)(\tilde x')-D(f\circ Q)(\tilde x)}\to0$.

\emph{Step 3 (induction).} We induct on $m$, for all normed $W$. The map $Q^*\colon\cL(P,W)\to\cL(P',W)$, $T_0\mapsto T_0\circ Q$, is linear with $\norm{Q^*T_0}\ge\norm{T_0}/\mathsf c$. By Step~2 its range is the set of operators that vanish on $\ker Q$, which is closed. By the chain rule, $D(f\circ Q)=Q^*\circ(Df)\circ Q$. If $f\circ Q$ is $C^{m+1}$, then $D(f\circ Q)$ is $C^m$ with values in the closed range of $Q^*$. By \cite[Lemma~2.4]{I} it is $C^m$ as a map into that range. Composing with the bounded inverse of $Q^*$ on its range, $(Df)\circ Q$ is $C^m$. By induction, applied to $Df\colon U\to\cL(P,W)$, $Df$ is $C^m$, so $f$ is $C^{m+1}$. For $m=\infty$ apply this for every finite $m$.
\end{proof}

So Theorem~\ref{dp:thm:minimal} reads: the DP germs are the smallest class that contains the analytic germs and has the two descents of $C^\infty$ maps.

\subsection{Ultrametric spaces}\label{dp:ssec:um-minimal}

Between ultrametric spaces the DP germs are already the AA germs, through an ultrametric enlargement \cite[Theorem~14.39(1)]{I}. Here ultrametric norms simplify further: no target descent is needed. The reason is in the bound of Lemma~\ref{dp:lem:laws-cover}. Its factor $n!$ counts the words that share a monomial, and removing it needs an enlargement of $W$. If $W$ is ultrametric and the cover has the supremum norm, the sum~\eqref{dp:eq:word-lift} has no count of terms and each $v_{\mathsf w}$ is bounded geometrically. So the cover alone makes each law a $W$-valued diagonal. In this subsection, unless said otherwise, the target of a germ is any normed space. Conditions (C1) and (C3) are read among the germs of each statement, so in (C3) the source $P'$ of $Q$ is also one of the sources of that statement.

\begin{theorem}[between ultrametric spaces]\label{dp:thm:ultrametric-minimal}
Let $P$ and $W$ be ultrametric normed spaces. Let $q^{\rm u}_P\colon c_{00}(I_P)\to P$, $e_h\mapsto h$, where $c_{00}(I_P)$ has the supremum norm. Then $q^{\rm u}_P$ is contractive with preimage constant $\abs\varpi^{-1}$. For $U\subseteq P$ open, $f\colon U\to W$ and $x_0=q^{\rm u}_P(\tilde x_0)\in U$:
\begin{enumerate}
\item $f$ is DP at $x_0$ if and only if $f\circ q^{\rm u}_P$ is analytic at $\tilde x_0$. Radius $R$ and bound $M$ give radius $\abs\varpi R$ and bound $M$, in both directions.
\item Among germs with ultrametric sources and ultrametric targets, the DP germs are the smallest class with (C1) and (C3).
\item The analytic germs between ultrametric spaces have (C3) if and only if every DP germ between ultrametric spaces is analytic. This fails for $K$ complete, not spherically complete, of characteristic $p$.
\item Let $K$ be complete and let $P$, $W$ be Banach. Items~(1) and~(2) hold, within Banach spaces, with the cover $c_0(I_P)\to P$, $\mathsf v\mapsto\sum_h\mathsf v_hh$.
\end{enumerate}
\end{theorem}

\begin{proof}
\emph{Step 1 (the cover).} As $P$ is ultrametric, $\norm{\sum_h\mathsf v_hh}\le\max_h\abs{\mathsf v_h}\norm h\le\norm{\mathsf v}_\infty$. The preimages are those of the paragraph after Definition~\ref{dp:def:cover}, of supremum norm $\abs\lambda^{-1}<\abs\varpi^{-1}\norm h$.

\emph{Step 2 (laws on the cover).} Let $J$ be a totally ordered set and $b\colon c_{00}(J)\to W$ a bounded law of degree $n$. Define $v_{\mathsf w}$ and $\beta_b$ as before~\eqref{dp:eq:word-lift}; the same computation gives $\beta_b(h,\dots,h)=b(h)$. On the span of the letters of $\mathsf w$, which carries the maximum norm, \cite[Lemma~Q.2(4)]{I} gives $\norm{v_{\mathsf w}}\le\abs\varpi^{-n}\norm b$. As $W$ is ultrametric, $\norm{\beta_b(h^1,\dots,h^n)}\le\max_{\mathsf w}\prod_c\abs{h^c_{\mathsf w_c}}\norm{v_{\mathsf w}}\le\abs\varpi^{-n}\norm b\prod_c\norm{h^c}_\infty$. So $\norm{\beta_b}\le\abs\varpi^{-n}\norm b$.

\emph{Step 3 (item (1)).} Let $f$ be DP at $x_0$ with radius $R$, bound $M$ and laws $a_n$. By Proposition~\ref{dp:prop:source-descent}(2), $f\circ q^{\rm u}_P$ is DP at $\tilde x_0$ with laws $a_n\circ q^{\rm u}_P$ of norm at most $MR^{-n}$. By Step~2 each is the diagonal of an $n$-linear map into $W$ of norm at most $M(\abs\varpi R)^{-n}$. So $f\circ q^{\rm u}_P$ is analytic at $\tilde x_0$ with radius $\abs\varpi R$ and bound $M$. Conversely, an analytic germ with radius $R$ and bound $M$ is DP with the same data \cite[Proposition~14.9(1)]{I}, and Proposition~\ref{dp:prop:source-descent}(1) with $\mathsf c=\abs\varpi^{-1}$ gives radius $\abs\varpi R$ and bound $M$.

\emph{Step 4 (item (2)).} The DP germs contain the analytic germs and have (C3). Let $\cC$ have (C1) and (C3), and let $f$ be DP at $x_0$. By item~(1), $f\circ q^{\rm u}_P$ is analytic, so it is in $\cC$. As $c_{00}(I_P)$ is ultrametric, the germ stays in the class of germs considered, and (C3) puts $f$ in $\cC$.

\emph{Step 5 (item (3)).} If the analytic germs have (C3) and $f$ is DP, then $f\circ q^{\rm u}_P$ is analytic by item~(1), so $f$ is analytic. Conversely, if the analytic germs are the DP germs, they have (C3) by Proposition~\ref{dp:prop:source-descent}(1). For the failure, \cite[Proposition~14.40(3)]{I} gives a function on an ultrametric Banach space that is AA, hence DP, at every point and analytic at no point.

\emph{Step 6 (item (4)).} The map $\mathsf v\mapsto\sum_h\mathsf v_hh$ is defined: the sum converges, as $\mathsf v_hh\to0$ in the complete ultrametric $P$. The map has the properties of Step~1. Let $b\colon c_0(J)\to W$ be a bounded law of degree $n$, with $v_{\mathsf w}$ as in Step~2. For $h^1,\dots,h^n\in c_0(J)$ and $\varepsilon>0$, only finitely many words have $\prod_c\abs{h^c_{\mathsf w_c}}>\varepsilon$. So~\eqref{dp:eq:word-lift} converges in the Banach space $W$, with $\norm{\beta_b}\le\abs\varpi^{-n}\norm b$. The maps $h\mapsto\beta_b(h,\dots,h)$ and $b$ are continuous (Theorem~\ref{dp:thm:calculus}(1)) and agree on the dense subspace $c_{00}(J)$. Steps~3 and~4 then apply.
\end{proof}

\begin{corollary}[ultrametric sources, every target]\label{dp:cor:um-cover}
Let $P$ be ultrametric and $W$ normed, and let $f$ be DP at $x_0$ with radius $R$ and bound $M$. Then $f\circ q_P$ is analytic at every preimage of $x_0$, with radius $R/C_1$ and bound $M$. So every class of germs that has (C1) for the sources $\ell^1_{00}(J)$ and (C3) for the canonical covers contains every DP germ on an ultrametric source, with no target descent.
\end{corollary}

\begin{proof}
Let $a_n$ be the laws of $f$ at $x_0$. By Lemma~\ref{dp:lem:lift-cover} and Lemma~\ref{dp:lem:strong-formal}(3), $a_n\circ q_P$ is the diagonal of a bounded $n$-linear map into $W$ of norm at most $\lvert\!\lvert\!\lvert a_n\rvert\!\rvert\!\rvert\le C_1^n\norm{a_n}\le M(R/C_1)^{-n}$. By Proposition~\ref{dp:prop:source-descent}(2), these laws expand $f\circ q_P$.
\end{proof}

For $P\ne0$ the source $\ell^1_{00}(I_P)$ of this cover is not ultrametric. In characteristic $0$ the DP germs on ultrametric sources are the analytic germs \cite[Proposition~14.40(2)]{I}, so (C1) alone gives them. In characteristic $p$, whether, among germs with ultrametric sources and Banach targets, the DP germs are the smallest class with (C1), (C2), (C3) is open; see Question~\ref{dp:q:c0}(2).

\subsection{What the descents reach}\label{dp:ssec:reach}

Let $(Z;W_1,W_2)$ and $(Z';W_1',W_2')$ be filtered spaces as in \cite[Definition~14.1]{I}, with $Z,Z'$ Banach, and with quotient maps $\operatorname{qt}\colon W_1\to W_1/W_2$ and $\operatorname{qt}'\colon W_1'\to W_1'/W_2'$. Let $w_0\in W_1$. A \emph{morphism germ} at $w_0$ is the germ of a map $\Phi$, from an open neighborhood $\mathsf O$ of $w_0$ in $Z$ into $Z'$, analytic at $w_0$, with $\Phi(\mathsf O\cap W_1)\subseteq W_1'$ and with $\Phi(w+w')-\Phi(w)\in W_2'$ whenever $w,w+w'\in\mathsf O\cap W_1$ and $w'\in W_2$. It induces a germ $\varphi$ at $\operatorname{qt}(w_0)$ from $W_1/W_2$ to $W_1'/W_2'$, by $\varphi(\operatorname{qt}(w))=\operatorname{qt}'(\Phi(w))$. As $Z'$ is complete, rearranging the series shows that a map analytic at $w_0$ is analytic on a ball around $w_0$, so this agrees with \cite[Definition~14.25]{I} when $W_2=W_2'=0$. If $W_2=0$, the induced germs are AA: $\Phi|_{W_1}$ is AA into the Banach space $W_1'$ \cite[Remark~14.13]{I}, and so is $\operatorname{qt}'\circ\Phi|_{W_1}$ \cite[Proposition~14.12(1)]{I}. With a quotient source, every DP germ between Banach spaces is induced, after the identifications of Proposition~\ref{dp:prop:subquotients}(2). For $K$ complete of characteristic $p$ these include germs that are AA at no point: the law $a$ of Remark~\ref{dp:rem:an-no-descent}, by Step~4 of the proof of Proposition~\ref{dp:prop:other-classes}, as AA at a point does not change under an equivalent norm on the source. In characteristic $0$ we do not decide whether they go beyond the AA germs: this is the case of Banach sources in \cite[Question~14.42(1)]{I}.

\begin{proposition}[subquotient sources]\label{dp:prop:subquotients}
\leavevmode
\begin{enumerate}
\item Every germ induced by a morphism germ between filtered spaces as above is DP at $\operatorname{qt}(w_0)$.
\item Let $K$ be complete, let $P,W$ be Banach spaces and let $f\colon U\to W$ be DP at $x_0\in U$. Then $f$ is induced by a morphism germ, after two identifications: of $P$, with an equivalent norm, with $\ell^1(I_P)/\ker\hat q_P$, and of $W$ with a closed subspace of a Banach space.
\end{enumerate}
\end{proposition}

\begin{proof}
\emph{Step 1 (item (1)).} As a map from $W_1$ into $Z'$, $\Phi|_{W_1}$ is analytic at $w_0$, hence DP at $w_0$ (Theorem~\ref{dp:thm:calculus}(3)). Its values near $w_0$ lie in the closed subspace $W_1'$, so it is DP into $W_1'$ (Theorem~\ref{dp:thm:target-descent}(2)). Composing with the bounded linear $\operatorname{qt}'$ keeps this. Near $w_0$ we have $\varphi\circ\operatorname{qt}=\operatorname{qt}'\circ\Phi|_{W_1}$, so $\varphi\circ\operatorname{qt}$ is DP at $w_0$. The map $\operatorname{qt}$ is open, so $\varphi$ is defined on the open set $\operatorname{qt}(\mathsf O\cap W_1)$, and $\operatorname{qt}$ has preimage constant $\abs\varpi^{-1}$ for the quotient norm. Proposition~\ref{dp:prop:source-descent}(1) gives that $\varphi$ is DP at $\operatorname{qt}(w_0)$.

\emph{Step 2 (item (2)).} By the proof of Theorem~\ref{dp:thm:minimal}(4), there are $\tilde x_0\in\ell^1(I_P)$ with $\hat q_P(\tilde x_0)=x_0$ and a linear isometry $j$ of $W$ into a Banach space $Z'$ with $\Phi=j\circ f\circ\hat q_P$ analytic at $\tilde x_0$. The map $\Phi$ has values in the closed subspace $j(W)$ and is constant on the cosets of $\ker\hat q_P$. So it is a morphism germ from $(\ell^1(I_P);\ell^1(I_P),\ker\hat q_P)$ to $(Z';j(W),0)$. As $\hat q_P$ is contractive with preimage constant $\abs\varpi^{-1}$, it induces a bijection of $\ell^1(I_P)/\ker\hat q_P$ onto $P$ with $\norm{\hat q_P(\mathsf v)}\le\norm{[\mathsf v]}\le\abs\varpi^{-1}\norm{\hat q_P(\mathsf v)}$. Under this bijection and $j$, the germ induced by $\Phi$ is $f$.
\end{proof}

An equivalent norm on the source changes the radius and keeps the class of DP germs. So, for $K$ complete and Banach spaces, the germs induced with quotient sources allowed are exactly the DP germs. This is Theorem~\ref{dp:thm:minimal} again: sub-pairs are target descent, and quotient sources are source descent.

Now the example that shows what this minimality rests on.

\begin{example}[the $\ell^1$ germ]\label{dp:ex:l1}
Let $P=\ell^1(\N,K)$ and $W=\Kh$. Choose $\mathsf e_n\in K$ with $n^n\le\abs{\mathsf e_n}<\abs\varpi^{-1}n^n$, and put
\[
 a_n(h)=\mathsf e_nh_1h_2\cdots h_n,\qquad\mathsf f(h)=\sum_{n\ge1}a_n(h)\qquad(\norm h<1).
\]
For $K$ complete this is the germ of \cite[Theorem~14.33(1)]{I}. Over every $K$:
\begin{enumerate}
\item $\norm{a_n}<\abs\varpi^{-1}$, so $\mathsf f$ is DP at $0$ with radius $1$ and bound $\abs\varpi^{-1}$.
\item $\mathsf f$ is AA at $0$ with radius $1/C_0$ and bound $\abs\varpi^{-1}$. So $j\circ\mathsf f$ is analytic at $0$ for a linear isometry $j$ of $\Kh$ into a Banach space.
\item $\mathsf f$ is not analytic at $0$.
\end{enumerate}
So the analytic germs fail (C2) in every characteristic. By items~(1) and~(3), finite coordinates cannot be dropped from Theorem~\ref{dp:thm:calculus}(3).
\end{example}

\begin{proof}
\emph{Step 1 (item (1)).} By the arithmetic--geometric mean inequality, $\prod_{i\le n}\abs{h_i}\le(\norm h/n)^n$. So $\norm{a_n}\le\abs{\mathsf e_n}n^{-n}<\abs\varpi^{-1}$.

\emph{Step 2 (item (2)).} By \cite[Theorem~14.39(2)]{I}, $\operatorname{amb}(a_n)\le C_0^n\norm{a_n}<\abs\varpi^{-1}C_0^n$. Then \cite[Proposition~14.9(1)]{I} and \cite[Remark~14.13]{I} apply.

\emph{Step 3 (item (3)).} Suppose $\mathsf f(h)=\sum_n\beta_n(h,\dots,h)$ near $0$, with $n$-linear $\beta_n\colon P^n\to\Kh$ and $\norm{\beta_n}\le MR^{-n}$. By Lemma~\ref{dp:lem:coefficients-unique}(2) with $j=\id$, $\beta_n(h,\dots,h)=a_n(h)$. As $K$ is infinite, comparing coefficients of $s_1\cdots s_n$ at $h=\sum_{i\le n}s_ie_i$ gives $\sum_{\sigma\in S_n}\beta_n(e_{\sigma(1)},\dots,e_{\sigma(n)})=\mathsf e_n$. The absolute value is nonarchimedean and $\norm{e_i}=1$, so $\norm{\beta_n}\ge\abs{\mathsf e_n}\ge n^n$. This is not at most $MR^{-n}$ for large $n$.
\end{proof}

\emph{The mechanism.} The norm of $\ell^1$ adds real numbers. So the arithmetic--geometric mean inequality applies, with a gain of $n^{-n}$, and $\mathsf e_nh_1\cdots h_n$ stays bounded with $\abs{\mathsf e_n}$ about $n^n$. A multilinear map $\beta$ with this diagonal has $\sum_\sigma\beta(e_{\sigma(1)},\dots,e_{\sigma(n)})=\mathsf e_n$. Over $\mathbb C$ the $n!$ orderings could share the load: $\beta$ could take values $\mathsf e_n/n!$, of size about $\mathrm e^n$. In $\Kh$ they cannot, as a sum equal to $\mathsf e_n$ has a term of absolute value at least $\abs{\mathsf e_n}$. Multilinear maps into an ultrametric space see only the supremum norm on $\ell^1$ (Proposition~\ref{dp:prop:hull}(1) and the paragraph before it), and there $h_1\cdots h_n$ has size $1$. So $\sum_na_n$ converges on the ball of $\ell^1$ because of an inequality between real numbers on the source. No divided power occurs: the example has the same form for $K=\mathbb C((t))$, in residue characteristic $0$. A non-ultrametric enlargement $Z$ can hold the shares in separate coordinates, which is why $j\circ\mathsf f$ is analytic for some $j$.

\emph{What this says about minimality.} By Example~\ref{dp:ex:l1}(2), every class with (C1) and (C2) contains $\mathsf f$. Among germs with Banach values, the class CA of germs analytic after a cover (Section~\ref{dp:sec:cover}) is the smallest class with (C1) and (C3), and it does not contain $\mathsf f$ (Lemma~\ref{dp:lem:cover-flat}(3), Theorem~\ref{dp:thm:cah}(5)). Within ultrametric targets, the strongly DP germs have (C1), (C3) and target descent along isometries between ultrametric spaces (Theorem~\ref{dp:thm:strong-descent}), and do not contain $\mathsf f$, as the coefficient of $a_n$ at $e_1,\dots,e_n$ has absolute value at least $n^n$. So the minimality of Theorem~\ref{dp:thm:minimal} on general sources rests on target descent into non-ultrametric spaces. On general sources, (C1) and (C3) alone do not reach the DP germs. Two axioms that reach them when added are (C2), by Theorem~\ref{dp:thm:minimal}, and the axiom (S) of Section~\ref{dp:sec:ledger}, by Proposition~\ref{dp:prop:s-axiom}(4). On ultrametric sources, (C1) and (C3) alone reach them, through covers that are not ultrametric (Corollary~\ref{dp:cor:um-cover}). Section~\ref{dp:sec:ledger} lists what each disagreement with analyticity is forced by.
